\scanpage{43}
enfin $\widetilde A\xrightarrow{\unsure{o_G|\sigma_G}}S$ est donné par passage au quotient à partir de $o_G|\widetilde B':\widetilde B'\to S'$, par commutativité dans

(51)\quad\diagram{\begin{tikzcd}
\widetilde B'\arrow[r,"o_{G'}|\widetilde B'"]\arrow[d,"\unsure{\sim|\wr}"']&S'=\text{sommets de }G'_\bullet\arrow[d]\\
\widetilde A\arrow[r]&S\simeq\pi_0(G'_\bullet)
\end{tikzcd}}

Ainsi $G$ est reconstruit à partir de $G'$, et $f$ de $\unsure{\overline m|=}$ bien sûr. On voit que $G$ apparaît comme un quotient de $G'$, et en fait c'est \unsure{plus|bien} une somme amalgamée dans la catégorie des graphes et des morphismes généralisés de graphes :

(52)\quad\diagram{\begin{tikzcd}[column sep=large,row sep=large]
G'_\bullet\arrow[r,hook,"\text{incl. can}"]\arrow[d,"\text{proj. can}"']\arrow[dr,phantom,"\text{cocart}"]&G'\arrow[d]\\
\begin{gathered}S=\pi_0(\unsure{G'_\bullet|S'_\bullet})\\
\text{\unsure{(graphe|Graphe} }\deleted{\unsure{\text{ponctuel}}}\inserted{\text{discret}}\\
\text{défini par }S\overset{\mathrm{dfn}}{=}\pi_0(G'_\bullet)
\end{gathered}\arrow[r]&G
\end{tikzcd}}

i.e. $G'\longrightarrow G$ est universel parmi les morphismes généralisés de \inserted{$G$ dans un} graphe $G'$, qui transforment les arcs de $\widetilde A'\setminus\widetilde B'$ en des identités. Posons donc
\mspage{42}
